\chapter{Il Deserto del Reale Digitale}
\label{cap:p4c39}

\epigraph{«Orwell temeva che la verità ci sarebbe stata nascosta. Huxley temeva che la verità sarebbe stata sommersa da un mare di irrilevanza.»}{Neil Postman, \textit{Divertirsi da morire} (1985)}

Il capitolo precedente si chiudeva con una danza che esce dalla stanza. Il giudizio che si adatta alla macchina è un fatto privato finché riguarda una persona alla volta; diventa un fatto pubblico quando ciò che produciamo insieme alle macchine si riversa, in quantità mai viste, nell'ambiente informativo che condividiamo. Quando Jean Baudrillard parlava del «deserto del reale», pensava alla realtà che scompare dietro il velo delle simulazioni. Non poteva immaginare che le simulazioni stesse si sarebbero moltiplicate fino a generarne altre, in un gioco di specchi che rischia di dimenticare che cosa stesse riflettendo.

Immaginate di entrare nella Biblioteca di Alessandria e scoprire che una parte crescente dei rotoli è stata vergata da automi instancabili, che rimescolano le parole dei veri studiosi in testi plausibili ma privi di quella scintilla che chiamiamo comprensione. Non è una descrizione letterale di Internet oggi, ma non è nemmeno pura fantascienza. Per questi contenuti si è diffuso un termine inglese, \textit{slop}, che indica materiale digitale di qualità mediocre generato in serie dall'intelligenza artificiale.

\section{Anatomia di un'inondazione}

Il nostro cervello, lo abbiamo visto nella Parte I, è molto più bravo a riconoscere schemi che a verificarli. Quando le informazioni sono troppe, la verifica è la prima cosa che salta, e a passare sono i contenuti più familiari, più popolari o più coinvolgenti, non i più accurati.\footnote{Menczer, F., e Hills, T. (2020). Information overload helps fake news spread, and social media knows it. \textit{Scientific American}, dicembre 2020. L'articolo precede la diffusione dell'IA generativa e riguarda il sovraccarico informativo in generale.} L'IA generativa non ha creato questa debolezza, ma abbassa quasi a zero il costo di sfruttarla.

Quella che con una certa ironia chiamiamo «AI slop» (traducibile come «brodaglia artificiale») non è solo una questione estetica. Assomiglia piuttosto all'introduzione di una specie invasiva in un ecosistema fragile. I dati sulla scala del fenomeno vanno letti con cautela. Secondo il rapporto annuale di Imperva, una società di sicurezza informatica che vende servizi di protezione dai bot, nel 2023 quasi metà del traffico web era generata da programmi automatici, e circa un terzo da bot con finalità dannose\footnote{Imperva (2024). \textit{Bad Bot Report 2024}. Si tratta di stime prodotte da un'azienda del settore, basate sul traffico osservato dai suoi clienti. Riguardano il traffico (visite, richieste automatiche, tentativi di frode), non la quota di contenuti generati dall'IA.}. Il traffico non è il contenuto, ma il quadro è quello di una rete in cui una parte notevole dell'attività non ha un essere umano dall'altra parte.

\index[main]{AI slop}
\index[cognitive]{riconoscimento di pattern}

La tecnologia ha reso la produzione di contenuti mediocri istantanea e quasi gratuita. Nelle librerie online compaiono libri generati automaticamente, guide di viaggio su luoghi che nessuno ha visitato, manuali assemblati da sistemi che non hanno mai applicato quelle conoscenze. Le piattaforme musicali ospitano brani generati in serie, e su YouTube prosperano video-saggi su scienza o \textit{true crime} con voce sintetica, copione scritto da un modello linguistico e miniature create da generatori di immagini.

Il prodotto finale è \textit{plausibile}, ed è proprio questo il punto critico. È abbastanza coerente da superare una verifica superficiale, abbastanza generico da risultare inoffensivo, ma privo della profondità che nasce dall'esperienza diretta. Un frutto dall'aspetto perfetto e dal sapore acquoso. La figura~\ref{fig:inondazione_slop} restituisce l'idea, non un'ondata di falsi evidenti ma una marea di quasi giusti.

\figurasola{images/inondazione_slop.png}{L'inondazione dell'AI slop: contenuti plausibili, abbondanti e a costo quasi nullo sommergono il poco che è curato.}{fig:inondazione_slop}

\section{Il paradosso della creatività}

La produzione automatizzata su larga scala poggia su fondamenta che sollevano questioni etiche profonde. I modelli che generano questi contenuti sono stati addestrati su un vasto corpus di creatività umana (video, discussioni, opere d'arte, articoli scientifici, libri nati da esperienze di vita), raccolto in gran parte con il web scraping, spesso in aree grigie quanto a consenso e attribuzione.

Il paradosso è quasi crudele. L'arte e il pensiero umano sono serviti da carburante per costruire sistemi che ora competono economicamente con gli stessi creatori. Viene in mente il mito di Crono che divora i propri figli.

C'è chi prende posizione. Il canale di divulgazione scientifica Kurzgesagt ha dichiarato pubblicamente che continuerà a privilegiare il processo creativo umano, usando l'IA per alcuni aspetti tecnici ma non come sostituto della ricerca e della scrittura\footnote{Kurzgesagt -- In a Nutshell, video dedicato all'AI slop pubblicato su YouTube nell'autunno 2025.}. La domanda che resta aperta è inquietante. Quanto può reggere questa scelta, quando produrre cento video mediocri costa meno che produrne uno eccellente?

\section{Una questione di fiducia}

Il punto critico, per la fiducia collettiva nell'informazione, non sono tanto gli errori evidenti quanto quella che potremmo chiamare \textit{quasi-correttezza}. Un modello linguistico può essere abbastanza preciso da apparire autorevole e intanto commettere errori con assoluta sicurezza, intessendo imprecisioni sottili in narrazioni che suonano plausibili\footnote{Bender, E. M., Gebru, T., McMillan-Major, A., \& Shmitchell, S. (2021). «On the Dangers of Stochastic Parrots: Can Language Models Be Too Big?». \textit{Proceedings of FAccT}, 610-623.}.

Somiglia a quello che fa il nostro archivista bugiardo (cap.~\ref{cap:p1c6}) quando un testimone in buona fede ricorda nei minimi dettagli eventi mai accaduti. Non mente, ricostruisce ciò che è plausibile. Un modello linguistico fa qualcosa di analogo per ragioni diverse, perché genera testo che \textit{suona} corretto avendo imparato i pattern statistici del linguaggio, non la verità sottostante. E, addestrato anche a produrre risposte gradite, può aggiungere i dettagli che rendono una storia più coerente e più attraente, senza intenzione e senza agenda.

\index[main]{disinformazione}
\index[main]{fiducia epistemica}

Già nel 2023 l'organizzazione NewsGuard, che valuta l'affidabilità dei siti di informazione, segnalava decine di siti di «notizie» che pubblicavano soprattutto contenuti generati dall'IA con scarsa o nessuna supervisione umana; nei suoi monitoraggi successivi il numero è cresciuto sensibilmente\footnote{NewsGuard (2023). «Rise of the Newsbots: AI-Generated News Websites Proliferating Online». NewsGuard, 1 maggio 2023. Il primo rapporto ne individuava 49; i conteggi successivi sono aggiornati periodicamente dalla stessa organizzazione.}. Questi testi vengono indicizzati, condivisi e possono finire nei dati con cui si addestrano altri modelli. Qui si chiude un cerchio che riguarda la domanda di questo libro. Le macchine imparano dalle tracce del passato; se una parte crescente di quelle tracce è stata prodotta da altre macchine, il futuro verrà costruito su copie di copie. Alcuni esperimenti suggeriscono che modelli addestrati ripetutamente su dati generati da modelli perdono via via varietà, a cominciare dai casi rari.\footnote{Shumailov, I., e altri (2024). AI models collapse when trained on recursively generated data. \textit{Nature}, 631, 755-759. Sono esperimenti in condizioni controllate; quanto il fenomeno pesi nei sistemi reali, che mescolano dati umani e sintetici, è ancora discusso.}

Persino la scrittura specialistica mostra segnali del cambiamento, come abbiamo visto nel cap.~\ref{cap:p3c27} con gli articoli biomedici. Un'analisi delle revisioni di articoli presentati ad alcune grandi conferenze di intelligenza artificiale ha osservato un aumento improvviso di aggettivi come «meticoloso» o «intricato», e stimato che una quota non trascurabile di quei testi fosse stata sostanzialmente modificata con un modello linguistico\footnote{Liang, W., e altri (2024). «Monitoring AI-Modified Content at Scale: A Case Study on the Impact of ChatGPT on AI Conference Peer Reviews». \textit{Proceedings of ICML 2024}; arXiv:2403.07183.}.

\section{Quantità, qualità, verità}
\label{sec:p4c39-tre-assi}

Nel parlare di «slop» si mescolano spesso tre problemi diversi, che conviene separare perché attivano bias diversi.

\begin{concetto}[Tre problemi da non confondere]
\textbf{Quantità.} Il costo di produzione di un contenuto si avvicina a zero, e i contenuti si moltiplicano. Il problema non è il singolo pezzo, ma il rapporto tra segnale e rumore: trovare ciò che vale diventa più faticoso.

\textbf{Qualità.} Molti contenuti generati in serie sono corretti ma generici: non sbagliano, semplicemente non aggiungono nulla. Un testo può essere vero e inutile.

\textbf{Verità.} Alcuni contenuti, generati o no, contengono affermazioni false. Qui il problema non è la mediocrità ma l'errore, involontario (un'«allucinazione» del modello) o voluto (il falso fabbricato apposta).
\end{concetto}

Le tre cose non vanno insieme per forza. Un video generato in serie può essere accurato; un saggio scritto a mano può essere falso; un contenuto eccellente può annegare nella quantità. La quantità sfrutta la nostra tendenza a fermarci alla prima cosa disponibile; la qualità mediocre passa inosservata grazie all'effetto alone della fluidità; la falsità approfitta del bias di verità illusoria, per cui ciò che abbiamo già visto più volte ci sembra più vero.\footnote{Hasher, L., Goldstein, D., e Toppino, T. (1977). Frequency and the conference of referential validity. \textit{Journal of Verbal Learning and Verbal Behavior}, 16(1), 107-112; Fazio, L. K., e altri (2015). Knowledge does not protect against illusory truth. \textit{Journal of Experimental Psychology: General}, 144(5), 993-1002.} L'ultimo meccanismo è il più insidioso, perché in un ambiente dove ripetere costa nulla la ripetizione stessa diventa una forma di prova.

\figurasola{images/quantita_qualita_verita.png}{Tre problemi da non confondere. Tanti contenuti non vuol dire contenuti mediocri, e un contenuto mediocre non è per forza falso. Ogni problema attiva bias diversi e chiede rimedi diversi.}{fig:quantita_qualita_verita}

Quanto ai formati brevi e frammentati in cui tutto questo viene spesso consumato, la tentazione di attribuire loro danni precisi al cervello è forte, ma le prove sono ancora in larga parte correlazionali.

\section{La ricerca di autenticità}

Come stanno rispondendo le persone? Alcuni osservatori ipotizzano che molti utenti possano ridurre gradualmente la propria presenza online. Quando un ambiente digitale sembra popolato soprattutto da bot che producono contenuti per altri bot, gli esseri umani potrebbero spostarsi altrove, come chi lascia una festa accorgendosi che gli invitati sono quasi tutti manichini ben vestiti. Altri, al contrario, usano i social sempre più come palcoscenico identitario, una socialità diversa, più simile a una vetrina che a un caffè con un amico.

\index[main]{autenticità}

In mezzo emerge una preferenza per l'autenticità verificabile. Molte persone mostrano una certa resistenza verso l'arte generata artificialmente, non per pregiudizio tecnologico, ma perché all'arte chiedono qualcosa che viene dal vissuto di un organismo che ha provato dolore, gioia, perdita. Ogni pennellata degli autoritratti di Frida Kahlo portava il peso della sua sofferenza fisica; un'intelligenza artificiale non ha mai sentito il sole sulla pelle né l'attesa ansiosa.

Per me questa domanda ha smesso da poco di essere astratta. Dopo quasi quindici anni da utente fedele di Spotify, l'algoritmo mi ha proposto un artista che non conoscevo, Ulrich Yande. La musica era notevole, la voce impressionante. Ma qualcosa non tornava, a cominciare da tre album pubblicati in un solo anno, il 2025. Ho provato a saperne di più, e le foto del suo profilo Instagram mi sono parse quelle di un'altra persona, con ogni probabilità rielaborate dall'IA. Non posso dimostrarlo, e non pretendo di farlo. Posso dire soltanto che non ho trovato traccia di una persona reale dietro quel profilo, e che il sospetto, una volta nato, non se n'è più andato. Del resto esistono ormai piattaforme (Suno è solo un esempio) che generano una canzone completa, voce e strumenti compresi, da una semplice descrizione testuale. Che il fenomeno esista, in ogni caso, è documentato. Nell'estate del 2025 una band chiamata The Velvet Sundown, arrivata in poche settimane a centinaia di migliaia di ascoltatori mensili su Spotify, ammise dopo molte smentite di essere un «progetto musicale sintetico», con brani, voci e immagini generati con l'IA.\footnote{Il caso fu ripreso dalla stampa internazionale nel luglio 2025, tra gli altri da \textit{The Guardian} e \textit{Rolling Stone}.} Nel mio caso, invece, resta un sospetto, e mi interessa la reazione che ha provocato in me. Per la prima volta sto seriamente pensando di disdire l'abbonamento. Non mi bastava che quella voce fosse bella, volevo sapere che dietro ci fosse qualcuno.

Il curioso è che il corpo, forse, la differenza la registra anche quando la mente non ci arriva. In uno studio del 2025 l'organismo degli ascoltatori ha reagito in modo sottilmente diverso alla musica generata dall'intelligenza artificiale rispetto a quella composta da umani, anche senza sapere quale stessero ascoltando.\footnote{Fi\v{s}er, N., Mart\'{i}n-Pascual, M. A., \& Andreu-S\'{a}nchez, C. (2025). «Emotional impact of AI-generated vs. human-composed music in audiovisual media: A biometric and self-report study». \textit{PLOS One}, 20(6), e0326498.} Con la musica sintetica le pupille si dilatavano di più, segno di maggiore sforzo cognitivo, mentre la risposta galvanica della pelle indicava un'attivazione più bassa. Eppure le emozioni riportate erano quasi identiche. È un singolo studio, su un campione piccolo, e va preso come un indizio, non come una prova.

Questo bisogno di sapere chi c'è dietro tocca anche le professioni. Chi si confida con il proprio terapeuta ha diritto di sapere se le risposte che riceve sono mediate da un algoritmo? Chi affronta un processo dovrebbe essere informato se il suo avvocato delega parti significative del ragionamento giuridico a un sistema che riorganizza pattern statistici senza comprenderne le implicazioni umane?

\section{L'arte della presenza}

Che cosa resta all'umano, in questo paesaggio? Forse non il fatto di essere sempre corretti (del resto nemmeno le macchine lo sono, per quanto suonino sicure), ma la possibilità di riconoscere i propri limiti. Quando diciamo «non lo so, ma posso cercare di capirlo insieme a te», facciamo qualcosa di diverso da un sistema che genera risposte con certezza statistica. Una conversazione faccia a faccia, in cui non possiamo cancellare ciò che abbiamo detto né prenderci dieci minuti per la risposta perfetta, conserva una qualità che la sovrabbondanza di contenuti levigati rende più rara, e forse più preziosa.

Non è un elogio dell'analogico né un rifiuto della tecnologia, che può sollevarci da compiti ripetitivi e rendere accessibili conoscenze un tempo riservate a pochi. È un'osservazione sul valore. In un ambiente che tende all'uniformità, ciò che è lento, imperfetto e riconoscibilmente umano diventa una forma di biodiversità. Il deserto del reale digitale si estende, ma i deserti, osservati con attenzione, non sono vuoti.

C'è però un limite a questa consolazione. La presenza e l'autenticità funzionano finché sappiamo riconoscerle, finché possiamo dire con ragionevole sicurezza che dietro una voce o un volto c'è davvero qualcuno. Il rumore di fondo dello slop logora questa capacità lentamente, senza che nessuno lo voglia. Ma lo stesso logoramento offre un'occasione a chi invece vuole ingannare. In un ambiente già saturo di quasi-vero, un falso fabbricato con cura non ha nemmeno bisogno di essere perfetto per passare. E quando a essere falsificati sono proprio il volto e la voce, cioè le prove su cui la nostra mente ha sempre fondato la fiducia, la fragilità della realtà condivisa smette di essere un effetto collaterale e diventa un bersaglio.
